\section{Resumen ejecutivo}

Este documento reúne y vuelve a interpretar toda la evidencia disponible del proyecto de multiplicación de matrices y regla del trapecio en un clúster de cuatro nodos con 96 núcleos físicos. Se mantienen separados los experimentos históricos, la batería ampliada y la campaña de optimización para que cada comparación conserve su protocolo y su alcance.

\begin{center}
\small
\begin{tabular}{p{3.5cm}r p{8.0cm}}
\toprule
Fuente & Registros & Qué permite afirmar \\
\midrule
Primera batería & 50 & Comparación inicial entre Ethernet y Wi-Fi con $P=1,24,48,72,96$; una corrida por configuración. \\
Batería ampliada & 75 & Repite tamaños y añade $N=512,4096,6144$; conserva TX+RX del servidor y comprobaciones de consistencia. \\
Optimización controlada & 118 & Medianas y rangos de 3 o 5 repeticiones; compara difusión jerárquica, SUMMA 2D, algoritmos tuned, tamaños mayores y reparto asimétrico. \\
\midrule
Total documentado & 243 & Registros de ejecución conservados en CSV, logs, metadatos y sellos SHA-256. \\
\bottomrule
\end{tabular}
\end{center}

\subsection{Respuestas breves}

\begin{center}
\small
\begin{tabularx}{\linewidth}{>{\raggedright\arraybackslash}p{4.3cm}X>{\raggedright\arraybackslash}p{2.6cm}}
\toprule
Pregunta & Respuesta sustentada por los datos & Evidencia \\
\midrule
¿Conviene usar los cuatro nodos para matrices? & No en ningún tamaño medido ($N\le7168$). El mejor tiempo siempre fue un nodo con 24 procesos; con 96 procesos, 1D global tardó entre 4.57 y 9.65 veces más por Ethernet. & Secc.~\ref{sec:opt}, tabla~\ref{tab:opt-resumen} \\
¿Qué optimización funcionó? & Difundir $B$ una vez por nodo. Con $N=3072$ y $P=96$, la mediana bajó de 6.260 a 3.217 s por Ethernet y de 101.082 a 51.479 s por Wi-Fi; el TX agregado bajó a menos de la mitad. & Secc.~\ref{sec:opt} \\
¿Hace falta programarla a mano? & En la difusión aislada, pedir a Open MPI el algoritmo \texttt{knomial} produjo un TX similar y una mediana 3\,\% menor que la versión jerárquica escrita en C. Falta medirlo en el programa completo. & Tabla~\ref{tab:opt-resumen} \\
¿Mejoró SUMMA 2D? & No. Con $P=64$ tardó 1.75 veces más que 1D por Ethernet y movió casi el doble de datos. & Tabla~\ref{tab:opt-resumen} \\
¿Conviene distribuir el trapecio? & Sí, con carga grande: $n=10^{11}$ alcanza un speedup de 68.56 con 96 procesos por Ethernet. Con $n=10^8$ por Wi-Fi, 96 procesos son más lentos que uno. & Secc.~\ref{sec:trap-grande} \\
¿Ayuda el reparto asimétrico P/E? & No de forma concluyente. Con 24 procesos y $n=10^{11}$, la mediana mejora 1.05 veces pero los rangos se solapan; con 96 procesos es 3\,\% más lento. & Tabla~\ref{tab:trap-reps} \\
¿Son estables las corridas únicas? & En las configuraciones de matrices comparables con 96 procesos, las corridas únicas quedaron a menos de un 2\,\% de las medianas posteriores. Con 24 procesos la campaña fue 19--37\,\% más lenta, sobre todo en el cálculo, y se compiló sin \texttt{-march=native}. & Secc.~\ref{sec:variabilidad} \\
\bottomrule
\end{tabularx}
\end{center}

\subsection{Qué cambió respecto de la imagen inicial}

La ejecución de madrugada con ocho procesos repartidos entre cuatro nodos (dos por nodo), $N=3072$ y Wi-Fi tardó 906.351 s y registró 889.744 s en el máximo de \texttt{MPI\_Bcast(B)}. Ese caso sigue siendo un antecedente válido, pero no se mezcla con las campañas nuevas: usó otro ejecutable, otro kernel y otra instrumentación. La campaña controlada repite la pregunta con afinidad fija, kernel por bloques, comprobaciones de huellas y una separación explícita entre 1D global y 1D jerárquica. Por eso ahora se puede atribuir la mejora a una hipótesis comprobada en esta topología, sin afirmar que el comando por sí solo revele el algoritmo interno de Open MPI.

\subsection{Qué añade esta versión revisada}

\begin{itemize}
\item Una \textbf{auditoría completa} de la evidencia (sección~\ref{sec:auditoria}): las 243 filas de resultados coinciden con sus logs, las huellas numéricas son consistentes y los sellos SHA-256 están intactos. También documenta los problemas de trazabilidad hallados en archivos auxiliares y borradores.
\item Una comparación entre las \textbf{corridas únicas y las medianas repetidas} (sección~\ref{sec:variabilidad}), que acota cuánto se puede confiar en cada cifra de las baterías.
\item La advertencia de que la campaña se compiló \textbf{sin} \texttt{-march=native -funroll-loops}, a diferencia de las baterías: sus comparaciones internas son válidas, pero sus tiempos absolutos de cálculo no son directamente comparables con los de las baterías.
\item Una tabla resumen de la campaña con cocientes, solapamiento de rangos y el resultado \texttt{knomial}, más la tabla de repeticiones del trapecio con el desbalance leído de cada log.
\item Conclusiones unificadas al final, después de toda la evidencia, y figuras de la campaña con valores anotados y en formato vectorial.
\end{itemize}

\subsection{Alcance y reproducibilidad}

Los tiempos provienen de \texttt{MPI\_Wtime}; los contadores de red se describen junto a su alcance y no se presentan como bytes útiles exactos de una colectiva. Cada corrida nueva conserva su comando, entorno, log y resultado en \texttt{resultados\_optimizacion/}. Los scripts de esta carpeta regeneran tablas, figuras, auditorías y este PDF sin depender de los nodos remotos; leen la evidencia del repositorio original sin modificarla.
